\begin{abstract}
% 242 woorden
Habsora (`The Gospel') is an AI-targeting system deployed by the Israel Defense Forces (IDF) that generates over one hundred strike recommendations per day, each reviewed by a commander in approximately twenty seconds. Although each recommendation is formally approved by a human operator, this compression of human judgement into a near-instantaneous approval means genuine oversight is largely absent. This paper applies \citet{info12090385}'s Comprehensive Human Oversight Framework (CHOF) to analyse Habsora across its technical, socio-technical, and governance layers. The CHOF treats human oversight not as a single requirement but as a structured set of conditions that can be satisfied or violated independently at each layer. Our analysis finds that Habsora fails across all three. We argue this is a structural design failure rather than operator intent, with consequences cascading across layers. At the technical layer, ensemble opacity prevents operators from understanding the basis of any recommendation. At the socio-technical layer, time pressure and absent explainability produce automation bias, reducing formal human authority to rubber-stamping. At the governance layer, the absence of pre-deployment review requirements and the diffusion of accountability undermine both IHL compliance and post-hoc redress. To address these failures we propose three mutually dependent controls: a contrastive explainability interface grounded in SHAP and Grad-CAM, a mandatory deliberation lock enforcing IHL-structured confirmation steps, and a Glass Box governance framework \cite{ijcai2019p802} enabling independent audit. Together, these restore the conditions for meaningful human control without requiring access to classified model internals.
\end{abstract}